\documentclass[conference]{IEEEtran}
\usepackage{cite}
\usepackage{amsmath,amssymb,amsfonts}
\usepackage{graphicx}

\begin{document}

\title{IRIS: Interactive Rendering \& Intelligent Spatial-modeling via Hybrid Neural Pipelines on Consumer-Grade GPUs}

\author{\IEEEauthorblockN{George David Tsitlauri}
\IEEEauthorblockA{\textit{Dept. of Informatics \& Telecommunications}\\
\textit{University of Thessaly, Greece}\\
gdtsitlauri@gmail.com}}

\maketitle

\begin{abstract}
This paper introduces IRIS, a real-time neural rendering and perception framework designed for commodity GPUs with a strict 4GB VRAM budget. On an NVIDIA GTX 1650, the committed IRIS artifacts sustain roughly 27--29 FPS rendering while keeping rendering memory near 0.1GB, reach 30.19 Mrays/s hash-encoding throughput, and maintain low-latency SLAM and agent-loop execution in a unified CUDA-first stack. The strongest claim of the current release is therefore not universal renderer superiority, but the practicality of a bounded neural rendering and perception pipeline on a constrained consumer GPU.
\end{abstract}

\section{Introduction}
Neural Radiance Fields (NeRF) have revolutionized 3D reconstruction but remain computationally expensive. IRIS addresses this by utilizing a hybrid C++/Python pipeline and FP16 optimizations.

\section{Methodology: PRISM-GL}
The PRISM-GL (Perceptual Real-time Iterative Scene Modeling) algorithm utilizes temporal reprojection and hash-encoding to minimize the computational footprint on Turing-architecture GPUs.

\section{Hash Encoding Implementation}
IRIS uses an Instant-NGP style multi-resolution hash encoder with $L=16$ levels, $F=2$ features per level, and a hash table size of $2^{19}$ entries per level. The per-level grid resolution grows geometrically from 16 to 512. Input 3D coordinates are quantized per level and mapped with spatial hashing to compact FP16 feature vectors, producing a 32-dimensional encoded representation. A compact MLP (64 hidden units, 2 layers) predicts RGB and density from encoded features and view direction.

\section{PRISM-GL Temporal Accumulation}
To suppress flickering under sparse ray sampling, IRIS applies temporal accumulation:
\begin{equation}
I_t^{acc} = (1 - \alpha) \cdot \hat{I}_{t-1 \rightarrow t} + \alpha \cdot I_t
\end{equation}
where $\hat{I}_{t-1 \rightarrow t}$ is the previous frame reprojected by motion vectors and $\alpha = 0.1$ in our experiments. This improves stability while preserving responsiveness for autonomous control.

\section{Autonomous Vision Results}
The autonomous agent consumes depth maps, flags near-field obstacles at $d_{min}<0.5m$, and executes a rule-based policy with states \{FORWARD, TURN\_LEFT, TURN\_RIGHT, STOP\}. Over 100 simulation steps, the agent maintains sub-5ms latency and produces consistent avoidance behavior while logging trajectory and decisions for downstream analysis.

\section{Results}
Experimental results were collected on an NVIDIA GTX 1650 (4GB VRAM), CUDA + PyTorch, FP16 execution.

\begin{table}[h]
\caption{Rendering Benchmarks (GTX 1650)}
\centering
\begin{tabular}{lccc}
\hline
Resolution & Precision & FPS & VRAM (MB) \\
\hline
720p (scale 0.125) & FP16 & 28.47 & 111 \\
720p (scale 0.125) & FP32 & 57.56 & 164 \\
1080p (scale 0.08) & FP16 & 27.81 & 103 \\
1080p (scale 0.08) & FP32 & 31.84 & 156 \\
\hline
\end{tabular}
\end{table}

\subsection{Interpretation Boundary}
The IRIS results are best read as repository-specific hardware evidence for a
consumer-GPU neural rendering and perception pipeline. They show that the
current implementation can maintain practical throughput, low memory footprint,
and fast perception latencies under the committed GTX 1650 setup. They are not
intended to prove a universal ordering between FP16 and FP32 across all neural
rendering workloads, nor to claim parity with much larger rendering or robotics
systems.

\begin{table}[h]
\caption{SLAM Accuracy and Speed}
\centering
\begin{tabular}{lcc}
\hline
Frames & Trajectory Error (m) & Latency / Frame (ms) \\
\hline
10 & 0.045 & 15.71 \\
50 & 0.245 & 9.06 \\
100 & 0.495 & 9.21 \\
\hline
\end{tabular}
\end{table}

\begin{table}[h]
\caption{Autonomous Agent Simulation}
\centering
\begin{tabular}{lcc}
\hline
Metric & Value & Notes \\
\hline
Mean latency & 1.54 ms & 100 steps \\
Obstacle avoid rate & 1.00 & 77 avoids, 0 collisions \\
Collisions & 0 & near-field hard stops \\
\hline
\end{tabular}
\end{table}

\subsection{Why the Current Release Still Stands Up Well}
The present IRIS release is stronger than a minimal prototype because it does
not stop at an isolated renderer benchmark. The repository combines rendering,
hash-encoding throughput, SLAM-style perception, and an autonomous agent loop
inside one reproducible hardware-constrained stack. This gives IRIS a credible
identity as an edge spatial-computing baseline even though its claims are
deliberately bounded to the committed consumer-GPU configuration.

\section{Conclusion}
IRIS demonstrates that a compact spatial-computing pipeline can run credibly on
an edge-class consumer GPU, providing a useful research baseline for
lightweight neural rendering, SLAM-style perception, and low-latency autonomous
decision loops.

\section{Target Venue Fit}
The presented system aligns with CVPR and SIGGRAPH Asia interests in neural rendering efficiency, real-time perception, and hardware-aware deployment.

\end{document}
